\documentclass[11pt,a4paper]{article}
\usepackage[margin=1in]{geometry}
\usepackage[hidelinks]{hyperref}
\usepackage{longtable}
\usepackage{array}
\usepackage{titlesec}
\usepackage{enumitem}
\usepackage{booktabs}
\usepackage{listings}
\usepackage[T1]{fontenc}
\setlist[itemize]{leftmargin=1.5em}
\setlist[enumerate]{leftmargin=1.8em}
\titleformat{\section}{\large\bfseries}{\thesection}{0.6em}{}
\titleformat{\subsection}{\normalsize\bfseries}{\thesubsection}{0.6em}{}

\lstset{
  basicstyle=\ttfamily\small,
  breaklines=true,
  breakatwhitespace=false,
  columns=fullflexible
}

\title{newdb Shared Bundle\\External Integration Manual (Expanded)}
\author{newdb Project}
\date{\today}

\begin{document}
\maketitle
\tableofcontents
\newpage

\section{Document Scope}
This manual is an expanded LaTeX edition aligned with
\texttt{newdb/scripts/examples/README\_shared\_bundle.md}. It covers:
\begin{itemize}
  \item Build and distribution of self-contained shared bundles.
  \item Full exported interface references (\texttt{newdb} and \texttt{gtest\_capi}).
  \item Command-driven external integration model.
  \item Unified result parsing and error handling strategy.
  \item Practical cross-language integration notes.
\end{itemize}

\section{Shared Bundle Delivery Model}
\subsection{Output directories}
\begin{itemize}
  \item VS/MSVC: \texttt{build/newdb-vs2026-mt-tests/shared\_bundle/Release}
  \item MinGW: \texttt{build/newdb-mingw/shared\_bundle/Release}
  \item WSL/Linux: \texttt{build/newdb-wsl-tests/shared\_bundle/Release}
\end{itemize}

\subsection{Build commands}
\begin{itemize}
  \item VS:
\begin{lstlisting}
cmake --build build/newdb-vs2026-mt-tests --config Release --target shared_bundle --parallel
\end{lstlisting}
  \item MinGW:
\begin{lstlisting}
cmake --build build/newdb-mingw --target shared_bundle --parallel
\end{lstlisting}
  \item WSL:
\begin{lstlisting}
cmake --build build/newdb-wsl-tests --target shared_bundle -j8
\end{lstlisting}
\end{itemize}

\subsection{Dependency precedence strategy}
\begin{itemize}
  \item Windows: \texttt{run\_*.cmd} launchers prioritize bundle-local DLL lookup.
  \item Linux/WSL: binaries are linked with \texttt{RPATH=\$ORIGIN}, prioritizing colocated \texttt{.so}.
  \item MinGW: runtime DLLs (\texttt{libstdc++-6.dll}, \texttt{libgcc\_s\_*.dll}, \texttt{libwinpthread-1.dll}) are bundled.
\end{itemize}

\section{newdb C API Full Reference}
\subsection{Header and constants}
Header: \texttt{newdb/engine/include/newdb/c\_api.h}
\begin{itemize}
  \item API version: \texttt{NEWDB\_C\_API\_VERSION\_MAJOR/MINOR/PATCH}
  \item ABI version: \texttt{NEWDB\_C\_API\_ABI\_VERSION}
\end{itemize}

\subsection{Error codes}
\begin{longtable}{>{\ttfamily}p{0.34\linewidth}p{0.58\linewidth}}
\toprule
Code & Meaning \\
\midrule
NEWDB\_OK & Success \\
NEWDB\_ERR\_INVALID\_ARGUMENT & Invalid argument or buffer contract violation \\
NEWDB\_ERR\_INVALID\_HANDLE & Invalid or already-destroyed session handle \\
NEWDB\_ERR\_EXECUTION\_FAILED & Command/business execution failed \\
NEWDB\_ERR\_INTERNAL & Internal engine failure \\
NEWDB\_ERR\_LOG\_IO & File/log I/O failure \\
NEWDB\_ERR\_SESSION\_TERMINATED & Session is terminated and cannot continue \\
\bottomrule
\end{longtable}

\subsection{Thread-safety contract}
\begin{itemize}
  \item Global metadata APIs are thread-safe.
  \item Per-session APIs are safe only under ``one handle per thread''.
  \item Sharing one handle across concurrent threads is not guaranteed safe.
\end{itemize}

\subsection{Function catalog}
\subsubsection*{Version and ABI}
\begin{itemize}
  \item \texttt{newdb\_version\_string}
  \item \texttt{newdb\_api\_version\_major/minor/patch}
  \item \texttt{newdb\_abi\_version}
  \item \texttt{newdb\_negotiate\_abi}
  \item \texttt{newdb\_error\_code\_string}
\end{itemize}

\subsubsection*{Utility}
\begin{itemize}
  \item \texttt{newdb\_sum}
  \item \texttt{newdb\_check\_schema\_file}
\end{itemize}

\subsubsection*{Session lifecycle}
\begin{itemize}
  \item \texttt{newdb\_session\_create}
  \item \texttt{newdb\_session\_destroy}
  \item \texttt{newdb\_session\_set\_table}
  \item \texttt{newdb\_session\_last\_error}
\end{itemize}

\subsubsection*{Execution and telemetry}
\begin{itemize}
  \item \texttt{newdb\_session\_execute}
  \item \texttt{newdb\_session\_runtime\_stats}
  \item \texttt{newdb\_session\_append\_runtime\_snapshot}
\end{itemize}

\subsection{Recommended lifecycle}
\begin{enumerate}
  \item Check ABI compatibility with \texttt{newdb\_negotiate\_abi}.
  \item Create session handle.
  \item Execute command stream via \texttt{newdb\_session\_execute}.
  \item Collect stats/snapshots as needed.
  \item Destroy session.
\end{enumerate}

\section{gtest\_capi Full Reference}
Header: \texttt{newdb/include/gtest\_capi.h}

\subsection{Capabilities}
\begin{itemize}
  \item Initialize and run GoogleTest from external language runtimes.
  \item Enumerate tests and fetch test metadata.
  \item Configure runtime options (filter/output/repeat/shuffle/seed/etc.).
  \item Read post-run counters and elapsed time.
\end{itemize}

\subsection{Main APIs}
\begin{itemize}
  \item Init and run:
  \texttt{gtest\_capi\_init\_from\_argv},
  \texttt{gtest\_capi\_run\_all}
  \item Listing and metadata:
  \texttt{gtest\_capi\_list\_tests},
  \texttt{gtest\_capi\_enumerate\_tests}
  \item Options:
  \texttt{set/get\_filter},
  \texttt{set/get\_output},
  \texttt{set/get\_repeat},
  \texttt{set/get\_shuffle},
  \texttt{set/get\_random\_seed},
  \texttt{set/get\_color},
  \texttt{set/get\_break\_on\_failure},
  \texttt{set/get\_throw\_on\_failure},
  \texttt{set/get\_catch\_exceptions},
  \texttt{set/get\_also\_run\_disabled\_tests},
  \texttt{set/get\_brief},
  \texttt{set/get\_print\_time}
  \item Generic option bridge:
  \texttt{gtest\_capi\_set\_option},
  \texttt{gtest\_capi\_get\_option}
  \item Counters:
  total suite/test, to-run count, passed/failed/skipped/disabled,
  reportable disabled, elapsed ms.
\end{itemize}

\section{Feature-Oriented External Integration}
From external systems, \texttt{newdb} exposes full engine capability through command execution:
\begin{itemize}
  \item DDL/schema management.
  \item DML and query paths.
  \item Transaction, conflict, and isolation controls.
  \item WAL/recovery controls.
  \item Runtime and maintenance controls (vacuum/group commit/hot index).
\end{itemize}

\subsection{Recommended wrapper layering}
\begin{itemize}
  \item DDL wrapper
  \item DML wrapper
  \item Query wrapper
  \item Txn wrapper
  \item Runtime tuning wrapper
\end{itemize}

\section{Command Templates}
\subsection{DDL}
\begin{lstlisting}
CREATE TABLE(users)
USE(users)
SET PRIMARY KEY(uid)
CREATE SCHEMA(finance_v1)
ALTER TABLE users SET SCHEMA(finance_v1)
RENAME TABLE(users_v2)
DROP TABLE(users_v2)
\end{lstlisting}

\subsection{DML and query}
\begin{lstlisting}
INSERT(1, alice, 1000, rd)
UPDATE(1, alice, 1200)
DELETE(1)
COUNT
WHERE(balance, >, 1000, AND, dept, =, rd)
PAGE(1, 20, id, desc)
SUM(balance)
AVG(balance)
\end{lstlisting}

\subsection{Transaction and tuning}
\begin{lstlisting}
BEGIN
SAVEPOINT(sp1)
ROLLBACK TO(sp1)
COMMIT
TXNISOLATION snapshot
WRITECONFLICT wait 2000
AUTOVACUUM threshold 500
WALSYNC normal 200
GROUPCOMMIT 10 128
\end{lstlisting}

\section{Result Parsing and Error Handling}
Always evaluate all three:
\begin{enumerate}
  \item \texttt{rc} from \texttt{newdb\_session\_execute}
  \item \texttt{newdb\_session\_last\_error}
  \item \texttt{output\_buf}
\end{enumerate}

\subsection{Recommended result model}
\begin{lstlisting}
ExecuteResult {
  rc: int
  last_error: int
  error_name: string
  output: string
  ok: bool
}
\end{lstlisting}

\subsection{Buffer sizing}
Use at least 4KB for simple commands and 16KB+ for complex query outputs.
Do not rely on text completeness when truncation is possible.

\section{Cross-language Practical Notes}
\subsection{Python/C\#/Node strategy}
\begin{itemize}
  \item Load main library via absolute path.
  \item Keep one session handle per worker/thread.
  \item Wrap command execution in one adapter method returning unified result object.
  \item Convert non-\texttt{NEWDB\_OK} into domain-specific exceptions.
\end{itemize}

\section{Operational Guidance}
\subsection{Telemetry}
Use:
\begin{itemize}
  \item \texttt{newdb\_session\_runtime\_stats} for real-time metric snapshots.
  \item \texttt{newdb\_session\_append\_runtime\_snapshot} for JSONL historical traces.
\end{itemize}

\subsection{Suggested key metrics}
\begin{itemize}
  \item Commit latency (p95/max)
  \item Write conflict and timeout counts
  \item WAL recovery elapsed and scanned records
  \item Vacuum trigger count and bytes reclaimed
\end{itemize}

\section{Distribution Checklist}
\begin{enumerate}
  \item Ship entire \texttt{shared\_bundle/Release} directory.
  \item Do not split individual DLL/SO from their runtime dependencies.
  \item Enforce absolute-path dynamic loading in host app.
  \item Validate ABI before first command execution.
  \item Persist command-level execution records.
\end{enumerate}

\section{Troubleshooting}
\begin{itemize}
  \item ``Module not found'': missing colocated dependencies.
  \item Unexpected library resolution: enforce absolute load and bundle launchers.
  \item WSL googletest download failures: use local \texttt{FETCHCONTENT\_SOURCE\_DIR\_GOOGLETEST}.
\end{itemize}

\end{document}

